% ------------------------------------------------------------------------
%  FMM implementation, group 5: near-field list grouping.
%  Include with % ------------------------------------------------------------------------
%  FMM implementation, group 5: near-field list grouping.
%  Include with % ------------------------------------------------------------------------
%  FMM implementation, group 5: near-field list grouping.
%  Include with % ------------------------------------------------------------------------
%  FMM implementation, group 5: near-field list grouping.
%  Include with \input{figures/fmm_g5}
%  Companion to the pipeline overview in figures/fmm_overview.tex.
% ------------------------------------------------------------------------
\input{figures/fmm_style}

\begin{figure}[!htb]
\centering
\begin{tikzpicture}[
    >={Latex[length=1.8mm]},
    g/.style={fmmgpu, text width=14.6cm},
    c/.style={fmmcpu, text width=14.6cm},
]
\node[g, anchor=north west] (c1) at (0,0) {%
  \fmmH{\texttt{fmm\_build\_p2p\_csr()}}
  \fmmK{p2p\_count\_kernel}{--- 1 thread / near pair $\rightarrow$ \texttt{cnt[A]}}
  \fmmK{cub::DeviceScan::ExclusiveSum}{--- $\rightarrow$ \texttt{d\_p2p\_off}}
  \fmmK{p2p\_scatter\_kernel}{--- 1 thread / near pair, atomic fill of \texttt{d\_p2p\_val}}};

\foreach \n in {c1}
  \node[fmmtag] at ([shift={(-1.4mm,-1.2mm)}]\n.north east) {GPU};

\end{tikzpicture}
\caption[Group 5: near-field list grouping]{\textbf{Group~5: near-field list grouping.} Solid boxes are
    GPU kernels, dashed boxes host steps; italics give the parallel mapping.}
\label{fig:fmm-g5}
\end{figure}

%  Companion to the pipeline overview in figures/fmm_overview.tex.
% ------------------------------------------------------------------------
% ------------------------------------------------------------------------
%  Shared styling for the FMM implementation figures:
%    fmm_overview.tex   -- the whole pipeline, group headlines only
%    fmm_g1..g8.tex     -- one detailed figure per group
%  Idempotent: every figure \input{figures/fmm_style} before its tikzpicture.
%  Needs tikz with the libraries positioning, fit and arrows.meta.
% ------------------------------------------------------------------------

% Entry macros. The mapping runs on from the kernel name in the same paragraph
% rather than on its own line -- a forced break per entry doubles the height.
\providecommand{\fmmH}[1]{{\footnotesize\bfseries #1}\par\vspace{0.5ex}}
\providecommand{\fmmK}[2]{{\ttfamily\scriptsize #1}\,{\tiny\itshape\color{black!62}#2}\par\vspace{0.35ex}}
\providecommand{\fmmN}[1]{{\tiny\color{black!62}#1}\par\vspace{0.35ex}}

% Legend, shared by the overview and by any detail figure that wants it.
\providecommand{\fmmlegend}{{\scriptsize
   \tikz[baseline=-0.5ex]\node[draw, line width=0.6pt, fill=white,
        minimum width=4mm, minimum height=2.6mm, inner sep=0pt]{};\;GPU kernel
   \quad
   \tikz[baseline=-0.5ex]\node[draw=black!55, line width=0.5pt,
        dash pattern=on 1.4pt off 1.2pt, fill=black!7,
        minimum width=4mm, minimum height=2.6mm, inner sep=0pt]{};\;host (CPU) step
   \quad {\itshape italics} $=$ parallel mapping
   \quad \texttt{cub::} $=$ CUB device-wide primitive (GPU)}}

\tikzset{
  fmmgpu/.style  ={draw, line width=0.6pt, fill=white, align=left, inner sep=1.5mm},
  fmmcpu/.style  ={fmmgpu, draw=black!55, dash pattern=on 1.4pt off 1.2pt,
                   fill=black!7},
  fmmnote/.style ={fmmgpu, draw=black!45, fill=black!4, rounded corners=2pt},
  fmmtag/.style  ={font=\tiny\sffamily, text=black!55, fill=white,
                   inner sep=0.5mm, anchor=north east},
  fmmflow/.style ={->, line width=0.7pt, black!70},
  fmmgrp/.style  ={draw=black!45, line width=0.9pt, rounded corners=2.5pt,
                   inner sep=2.4mm},
  fmmlbl/.style  ={font=\small, fill=white, inner sep=1mm, anchor=west},
  fmmedge/.style ={font=\tiny\sffamily, text=black!60, inner sep=0.7mm,
                   fill=white},
  fmmband/.style ={font=\scriptsize\itshape, text=black!60, inner sep=1mm,
                   fill=white},
}


\begin{figure}[!htb]
\centering
\begin{tikzpicture}[
    >={Latex[length=1.8mm]},
    g/.style={fmmgpu, text width=14.6cm},
    c/.style={fmmcpu, text width=14.6cm},
]
\node[g, anchor=north west] (c1) at (0,0) {%
  \fmmH{\texttt{fmm\_build\_p2p\_csr()}}
  \fmmK{p2p\_count\_kernel}{--- 1 thread / near pair $\rightarrow$ \texttt{cnt[A]}}
  \fmmK{cub::DeviceScan::ExclusiveSum}{--- $\rightarrow$ \texttt{d\_p2p\_off}}
  \fmmK{p2p\_scatter\_kernel}{--- 1 thread / near pair, atomic fill of \texttt{d\_p2p\_val}}};

\foreach \n in {c1}
  \node[fmmtag] at ([shift={(-1.4mm,-1.2mm)}]\n.north east) {GPU};

\end{tikzpicture}
\caption[Group 5: near-field list grouping]{\textbf{Group~5: near-field list grouping.} Solid boxes are
    GPU kernels, dashed boxes host steps; italics give the parallel mapping.}
\label{fig:fmm-g5}
\end{figure}

%  Companion to the pipeline overview in figures/fmm_overview.tex.
% ------------------------------------------------------------------------
% ------------------------------------------------------------------------
%  Shared styling for the FMM implementation figures:
%    fmm_overview.tex   -- the whole pipeline, group headlines only
%    fmm_g1..g8.tex     -- one detailed figure per group
%  Idempotent: every figure % ------------------------------------------------------------------------
%  Shared styling for the FMM implementation figures:
%    fmm_overview.tex   -- the whole pipeline, group headlines only
%    fmm_g1..g8.tex     -- one detailed figure per group
%  Idempotent: every figure \input{figures/fmm_style} before its tikzpicture.
%  Needs tikz with the libraries positioning, fit and arrows.meta.
% ------------------------------------------------------------------------

% Entry macros. The mapping runs on from the kernel name in the same paragraph
% rather than on its own line -- a forced break per entry doubles the height.
\providecommand{\fmmH}[1]{{\footnotesize\bfseries #1}\par\vspace{0.5ex}}
\providecommand{\fmmK}[2]{{\ttfamily\scriptsize #1}\,{\tiny\itshape\color{black!62}#2}\par\vspace{0.35ex}}
\providecommand{\fmmN}[1]{{\tiny\color{black!62}#1}\par\vspace{0.35ex}}

% Legend, shared by the overview and by any detail figure that wants it.
\providecommand{\fmmlegend}{{\scriptsize
   \tikz[baseline=-0.5ex]\node[draw, line width=0.6pt, fill=white,
        minimum width=4mm, minimum height=2.6mm, inner sep=0pt]{};\;GPU kernel
   \quad
   \tikz[baseline=-0.5ex]\node[draw=black!55, line width=0.5pt,
        dash pattern=on 1.4pt off 1.2pt, fill=black!7,
        minimum width=4mm, minimum height=2.6mm, inner sep=0pt]{};\;host (CPU) step
   \quad {\itshape italics} $=$ parallel mapping
   \quad \texttt{cub::} $=$ CUB device-wide primitive (GPU)}}

\tikzset{
  fmmgpu/.style  ={draw, line width=0.6pt, fill=white, align=left, inner sep=1.5mm},
  fmmcpu/.style  ={fmmgpu, draw=black!55, dash pattern=on 1.4pt off 1.2pt,
                   fill=black!7},
  fmmnote/.style ={fmmgpu, draw=black!45, fill=black!4, rounded corners=2pt},
  fmmtag/.style  ={font=\tiny\sffamily, text=black!55, fill=white,
                   inner sep=0.5mm, anchor=north east},
  fmmflow/.style ={->, line width=0.7pt, black!70},
  fmmgrp/.style  ={draw=black!45, line width=0.9pt, rounded corners=2.5pt,
                   inner sep=2.4mm},
  fmmlbl/.style  ={font=\small, fill=white, inner sep=1mm, anchor=west},
  fmmedge/.style ={font=\tiny\sffamily, text=black!60, inner sep=0.7mm,
                   fill=white},
  fmmband/.style ={font=\scriptsize\itshape, text=black!60, inner sep=1mm,
                   fill=white},
}
 before its tikzpicture.
%  Needs tikz with the libraries positioning, fit and arrows.meta.
% ------------------------------------------------------------------------

% Entry macros. The mapping runs on from the kernel name in the same paragraph
% rather than on its own line -- a forced break per entry doubles the height.
\providecommand{\fmmH}[1]{{\footnotesize\bfseries #1}\par\vspace{0.5ex}}
\providecommand{\fmmK}[2]{{\ttfamily\scriptsize #1}\,{\tiny\itshape\color{black!62}#2}\par\vspace{0.35ex}}
\providecommand{\fmmN}[1]{{\tiny\color{black!62}#1}\par\vspace{0.35ex}}

% Legend, shared by the overview and by any detail figure that wants it.
\providecommand{\fmmlegend}{{\scriptsize
   \tikz[baseline=-0.5ex]\node[draw, line width=0.6pt, fill=white,
        minimum width=4mm, minimum height=2.6mm, inner sep=0pt]{};\;GPU kernel
   \quad
   \tikz[baseline=-0.5ex]\node[draw=black!55, line width=0.5pt,
        dash pattern=on 1.4pt off 1.2pt, fill=black!7,
        minimum width=4mm, minimum height=2.6mm, inner sep=0pt]{};\;host (CPU) step
   \quad {\itshape italics} $=$ parallel mapping
   \quad \texttt{cub::} $=$ CUB device-wide primitive (GPU)}}

\tikzset{
  fmmgpu/.style  ={draw, line width=0.6pt, fill=white, align=left, inner sep=1.5mm},
  fmmcpu/.style  ={fmmgpu, draw=black!55, dash pattern=on 1.4pt off 1.2pt,
                   fill=black!7},
  fmmnote/.style ={fmmgpu, draw=black!45, fill=black!4, rounded corners=2pt},
  fmmtag/.style  ={font=\tiny\sffamily, text=black!55, fill=white,
                   inner sep=0.5mm, anchor=north east},
  fmmflow/.style ={->, line width=0.7pt, black!70},
  fmmgrp/.style  ={draw=black!45, line width=0.9pt, rounded corners=2.5pt,
                   inner sep=2.4mm},
  fmmlbl/.style  ={font=\small, fill=white, inner sep=1mm, anchor=west},
  fmmedge/.style ={font=\tiny\sffamily, text=black!60, inner sep=0.7mm,
                   fill=white},
  fmmband/.style ={font=\scriptsize\itshape, text=black!60, inner sep=1mm,
                   fill=white},
}


\begin{figure}[!htb]
\centering
\begin{tikzpicture}[
    >={Latex[length=1.8mm]},
    g/.style={fmmgpu, text width=14.6cm},
    c/.style={fmmcpu, text width=14.6cm},
]
\node[g, anchor=north west] (c1) at (0,0) {%
  \fmmH{\texttt{fmm\_build\_p2p\_csr()}}
  \fmmK{p2p\_count\_kernel}{--- 1 thread / near pair $\rightarrow$ \texttt{cnt[A]}}
  \fmmK{cub::DeviceScan::ExclusiveSum}{--- $\rightarrow$ \texttt{d\_p2p\_off}}
  \fmmK{p2p\_scatter\_kernel}{--- 1 thread / near pair, atomic fill of \texttt{d\_p2p\_val}}};

\foreach \n in {c1}
  \node[fmmtag] at ([shift={(-1.4mm,-1.2mm)}]\n.north east) {GPU};

\end{tikzpicture}
\caption[Group 5: near-field list grouping]{\textbf{Group~5: near-field list grouping.} Solid boxes are
    GPU kernels, dashed boxes host steps; italics give the parallel mapping.}
\label{fig:fmm-g5}
\end{figure}

%  Companion to the pipeline overview in figures/fmm_overview.tex.
% ------------------------------------------------------------------------
% ------------------------------------------------------------------------
%  Shared styling for the FMM implementation figures:
%    fmm_overview.tex   -- the whole pipeline, group headlines only
%    fmm_g1..g8.tex     -- one detailed figure per group
%  Idempotent: every figure % ------------------------------------------------------------------------
%  Shared styling for the FMM implementation figures:
%    fmm_overview.tex   -- the whole pipeline, group headlines only
%    fmm_g1..g8.tex     -- one detailed figure per group
%  Idempotent: every figure % ------------------------------------------------------------------------
%  Shared styling for the FMM implementation figures:
%    fmm_overview.tex   -- the whole pipeline, group headlines only
%    fmm_g1..g8.tex     -- one detailed figure per group
%  Idempotent: every figure \input{figures/fmm_style} before its tikzpicture.
%  Needs tikz with the libraries positioning, fit and arrows.meta.
% ------------------------------------------------------------------------

% Entry macros. The mapping runs on from the kernel name in the same paragraph
% rather than on its own line -- a forced break per entry doubles the height.
\providecommand{\fmmH}[1]{{\footnotesize\bfseries #1}\par\vspace{0.5ex}}
\providecommand{\fmmK}[2]{{\ttfamily\scriptsize #1}\,{\tiny\itshape\color{black!62}#2}\par\vspace{0.35ex}}
\providecommand{\fmmN}[1]{{\tiny\color{black!62}#1}\par\vspace{0.35ex}}

% Legend, shared by the overview and by any detail figure that wants it.
\providecommand{\fmmlegend}{{\scriptsize
   \tikz[baseline=-0.5ex]\node[draw, line width=0.6pt, fill=white,
        minimum width=4mm, minimum height=2.6mm, inner sep=0pt]{};\;GPU kernel
   \quad
   \tikz[baseline=-0.5ex]\node[draw=black!55, line width=0.5pt,
        dash pattern=on 1.4pt off 1.2pt, fill=black!7,
        minimum width=4mm, minimum height=2.6mm, inner sep=0pt]{};\;host (CPU) step
   \quad {\itshape italics} $=$ parallel mapping
   \quad \texttt{cub::} $=$ CUB device-wide primitive (GPU)}}

\tikzset{
  fmmgpu/.style  ={draw, line width=0.6pt, fill=white, align=left, inner sep=1.5mm},
  fmmcpu/.style  ={fmmgpu, draw=black!55, dash pattern=on 1.4pt off 1.2pt,
                   fill=black!7},
  fmmnote/.style ={fmmgpu, draw=black!45, fill=black!4, rounded corners=2pt},
  fmmtag/.style  ={font=\tiny\sffamily, text=black!55, fill=white,
                   inner sep=0.5mm, anchor=north east},
  fmmflow/.style ={->, line width=0.7pt, black!70},
  fmmgrp/.style  ={draw=black!45, line width=0.9pt, rounded corners=2.5pt,
                   inner sep=2.4mm},
  fmmlbl/.style  ={font=\small, fill=white, inner sep=1mm, anchor=west},
  fmmedge/.style ={font=\tiny\sffamily, text=black!60, inner sep=0.7mm,
                   fill=white},
  fmmband/.style ={font=\scriptsize\itshape, text=black!60, inner sep=1mm,
                   fill=white},
}
 before its tikzpicture.
%  Needs tikz with the libraries positioning, fit and arrows.meta.
% ------------------------------------------------------------------------

% Entry macros. The mapping runs on from the kernel name in the same paragraph
% rather than on its own line -- a forced break per entry doubles the height.
\providecommand{\fmmH}[1]{{\footnotesize\bfseries #1}\par\vspace{0.5ex}}
\providecommand{\fmmK}[2]{{\ttfamily\scriptsize #1}\,{\tiny\itshape\color{black!62}#2}\par\vspace{0.35ex}}
\providecommand{\fmmN}[1]{{\tiny\color{black!62}#1}\par\vspace{0.35ex}}

% Legend, shared by the overview and by any detail figure that wants it.
\providecommand{\fmmlegend}{{\scriptsize
   \tikz[baseline=-0.5ex]\node[draw, line width=0.6pt, fill=white,
        minimum width=4mm, minimum height=2.6mm, inner sep=0pt]{};\;GPU kernel
   \quad
   \tikz[baseline=-0.5ex]\node[draw=black!55, line width=0.5pt,
        dash pattern=on 1.4pt off 1.2pt, fill=black!7,
        minimum width=4mm, minimum height=2.6mm, inner sep=0pt]{};\;host (CPU) step
   \quad {\itshape italics} $=$ parallel mapping
   \quad \texttt{cub::} $=$ CUB device-wide primitive (GPU)}}

\tikzset{
  fmmgpu/.style  ={draw, line width=0.6pt, fill=white, align=left, inner sep=1.5mm},
  fmmcpu/.style  ={fmmgpu, draw=black!55, dash pattern=on 1.4pt off 1.2pt,
                   fill=black!7},
  fmmnote/.style ={fmmgpu, draw=black!45, fill=black!4, rounded corners=2pt},
  fmmtag/.style  ={font=\tiny\sffamily, text=black!55, fill=white,
                   inner sep=0.5mm, anchor=north east},
  fmmflow/.style ={->, line width=0.7pt, black!70},
  fmmgrp/.style  ={draw=black!45, line width=0.9pt, rounded corners=2.5pt,
                   inner sep=2.4mm},
  fmmlbl/.style  ={font=\small, fill=white, inner sep=1mm, anchor=west},
  fmmedge/.style ={font=\tiny\sffamily, text=black!60, inner sep=0.7mm,
                   fill=white},
  fmmband/.style ={font=\scriptsize\itshape, text=black!60, inner sep=1mm,
                   fill=white},
}
 before its tikzpicture.
%  Needs tikz with the libraries positioning, fit and arrows.meta.
% ------------------------------------------------------------------------

% Entry macros. The mapping runs on from the kernel name in the same paragraph
% rather than on its own line -- a forced break per entry doubles the height.
\providecommand{\fmmH}[1]{{\footnotesize\bfseries #1}\par\vspace{0.5ex}}
\providecommand{\fmmK}[2]{{\ttfamily\scriptsize #1}\,{\tiny\itshape\color{black!62}#2}\par\vspace{0.35ex}}
\providecommand{\fmmN}[1]{{\tiny\color{black!62}#1}\par\vspace{0.35ex}}

% Legend, shared by the overview and by any detail figure that wants it.
\providecommand{\fmmlegend}{{\scriptsize
   \tikz[baseline=-0.5ex]\node[draw, line width=0.6pt, fill=white,
        minimum width=4mm, minimum height=2.6mm, inner sep=0pt]{};\;GPU kernel
   \quad
   \tikz[baseline=-0.5ex]\node[draw=black!55, line width=0.5pt,
        dash pattern=on 1.4pt off 1.2pt, fill=black!7,
        minimum width=4mm, minimum height=2.6mm, inner sep=0pt]{};\;host (CPU) step
   \quad {\itshape italics} $=$ parallel mapping
   \quad \texttt{cub::} $=$ CUB device-wide primitive (GPU)}}

\tikzset{
  fmmgpu/.style  ={draw, line width=0.6pt, fill=white, align=left, inner sep=1.5mm},
  fmmcpu/.style  ={fmmgpu, draw=black!55, dash pattern=on 1.4pt off 1.2pt,
                   fill=black!7},
  fmmnote/.style ={fmmgpu, draw=black!45, fill=black!4, rounded corners=2pt},
  fmmtag/.style  ={font=\tiny\sffamily, text=black!55, fill=white,
                   inner sep=0.5mm, anchor=north east},
  fmmflow/.style ={->, line width=0.7pt, black!70},
  fmmgrp/.style  ={draw=black!45, line width=0.9pt, rounded corners=2.5pt,
                   inner sep=2.4mm},
  fmmlbl/.style  ={font=\small, fill=white, inner sep=1mm, anchor=west},
  fmmedge/.style ={font=\tiny\sffamily, text=black!60, inner sep=0.7mm,
                   fill=white},
  fmmband/.style ={font=\scriptsize\itshape, text=black!60, inner sep=1mm,
                   fill=white},
}


\begin{figure}[!htb]
\centering
\begin{tikzpicture}[
    >={Latex[length=1.8mm]},
    g/.style={fmmgpu, text width=14.6cm},
    c/.style={fmmcpu, text width=14.6cm},
]
\node[g, anchor=north west] (c1) at (0,0) {%
  \fmmH{\texttt{fmm\_build\_p2p\_csr()}}
  \fmmK{p2p\_count\_kernel}{--- 1 thread / near pair $\rightarrow$ \texttt{cnt[A]}}
  \fmmK{cub::DeviceScan::ExclusiveSum}{--- $\rightarrow$ \texttt{d\_p2p\_off}}
  \fmmK{p2p\_scatter\_kernel}{--- 1 thread / near pair, atomic fill of \texttt{d\_p2p\_val}}};

\foreach \n in {c1}
  \node[fmmtag] at ([shift={(-1.4mm,-1.2mm)}]\n.north east) {GPU};

\end{tikzpicture}
\caption[Group 5: near-field list grouping]{\textbf{Group~5: near-field list grouping.} Solid boxes are
    GPU kernels, dashed boxes host steps; italics give the parallel mapping.}
\label{fig:fmm-g5}
\end{figure}
